\chapter{Introducción}
\label{cap:introduccion}

En este capítulo se presenta el contexto del proyecto, centrado en la escalada en
rocódromo y en el uso de técnicas de gamificación aplicadas al deporte. A partir de
este marco, se justifica la motivación del trabajo, se definen sus objetivos y se
resume la estructura de la memoria.


\section{Contexto de la escalada en rocódromo}

\subsection{Cómo funciona la escalada en rocódromo}

Un rocódromo es una instalación deportiva que simula la escalada en roca natural
mediante muros artificiales equipados con presas (agarres) de diferentes formas,
tamaños y colores. El objetivo del escalador es ascender por una ruta predefinida
utilizando exclusivamente las presas asignadas a dicha ruta.

Las rutas se identifican visualmente, habitualmente por el color de las presas. Cada
ruta tiene un punto de inicio marcado y un punto final o \textit{top} (última presa
que se debe controlar con ambas manos para dar el intento por válido). Cuando una
ruta se completa en el primer intento, se denomina \textit{flash}.

\subsection{Modalidades de escalada en rocódromo}

Las principales modalidades practicadas en rocódromo son la escalada de bloque
(\textit{búlder}) y la escalada deportiva (vía):

\begin{enumerate}
	\item \textbf{Escalada de bloque (búlder).} Se practica en muros de baja altura,
	normalmente de hasta 4 o 5 metros. El objetivo es resolver problemas o secuencias
	de movimientos cortos e intensos. Para la seguridad se emplean colchonetas
	(\textit{crash pads}) que amortiguan las caídas.

	\item \textbf{Escalada deportiva (vía).} Se realiza en muros de mayor altura y
	requiere completar recorridos largos donde predominan la resistencia y la gestión
	del esfuerzo. Es necesario usar cuerda, arnés y sistema de aseguramiento. Durante
	el ascenso, el escalador chapa la cuerda en anclajes intermedios y el asegurador,
	desde el suelo, controla la cuerda para detener posibles caídas.
\end{enumerate}

\subsection{Dificultad y progreso}

Cada ruta, tanto de bloque como de vía, tiene asignado un grado de dificultad que
permite estimar el nivel exigido, seleccionar desafíos adecuados y hacer seguimiento
del progreso. Este grado suele aparecer en la etiqueta de inicio.

En bloque se utilizan, entre otras, la escala de Fontainebleau (por ejemplo, 6A,
7A) y la \textit{V-Scale} (por ejemplo, V3, V6). En vía son frecuentes la escala
francesa (por ejemplo, 6a+, 7b) y la YDS (por ejemplo, 5.10a, 5.11c). Una explicación
detallada de estas escalas y sus conversiones puede consultarse en\footnote{\url{https://www.lacrux.com/es/escalar/Escalas-de-escalada-explicadas-UIAA-Fontainebleau-V-Grado-CO/}}.

Además, muchos rocódromos emplean sistemas internos basados en colores para asociar
rangos de dificultad a cada ruta. Estos sistemas facilitan la lectura del nivel y
pueden actuar como elemento motivador al permitir visualizar avances. Existen también
análisis estadísticos de ascensos que aportan referencias de progresión, aunque no
siempre son totalmente representativos de la población general de escaladores\footnote{\url{https://www.alessandromasullo.com/blog/analysis-of-4-million-climbing-ascents/}}.


\section{Contexto de la gamificación en deportes}

La gamificación consiste en aplicar principios y elementos propios de los juegos en
contextos no lúdicos, como el entrenamiento deportivo. Su finalidad es incrementar la
motivación, mejorar la retención de usuarios y favorecer la adherencia a la práctica
deportiva.

\subsection{Psicología: modelo RAMP}

Una base habitual para diseñar experiencias gamificadas es el modelo RAMP
(\textit{Relatedness, Autonomy, Mastery, Purpose}), alineado con teorías de
motivación intrínseca. Este modelo destaca cuatro necesidades principales:

\begin{itemize}
	\item \textbf{Relatedness (vinculación):} necesidad de conexión social,
	reconocimiento y pertenencia a una comunidad.
	\item \textbf{Autonomy (autonomía):} necesidad de control sobre las decisiones y
	la forma de progresar.
	\item \textbf{Mastery (maestría):} necesidad de mejora continua mediante retos
	alcanzables y progresivos.
	\item \textbf{Purpose (finalidad):} necesidad de que el esfuerzo tenga un
	significado más allá de la tarea inmediata.
\end{itemize}

Más información sobre motivación y gamificación puede consultarse en\footnote{\url{https://playmotiv.com/gamificacion-y-motivacion/}}.

\subsection{Sistema PBL (\textit{Points, Badges, Leaderboards})}

El sistema PBL es una implementación clásica de gamificación y se compone de tres
elementos:

\begin{itemize}
	\item \textbf{Points (puntos):} recompensan el esfuerzo o rendimiento del usuario
	(por ejemplo, dificultad encadenada, número de intentos o constancia).
	\item \textbf{Badges (insignias):} representan logros alcanzados y visualizan la
	evolución del deportista.
	\item \textbf{Leaderboards (clasificaciones):} comparan resultados entre usuarios,
	incentivando una competitividad sana y el sentido de comunidad.
\end{itemize}

Una introducción al modelo PBL se describe en\footnote{\url{https://www.adrformacion.com/knowledge/educacion-y-formacion/la_gamificacion_y_la_triada_pbl.html}}.

\subsection{Beneficios de la gamificación deportiva}

La integración de mecánicas gamificadas en aplicaciones deportivas aporta beneficios
relevantes:

\begin{itemize}
	\item \textbf{Mayor retención:} la retroalimentación constante reduce el abandono
	y mejora la adherencia al entrenamiento.
	\item \textbf{Visualización del progreso:} el registro de datos permite observar
	evolución objetiva mediante métricas y logros.
	\item \textbf{Refuerzo de comunidad:} la interacción social y la comparación
	amistosa aumentan la implicación de los usuarios.
\end{itemize}

En un deporte como la escalada, donde el progreso puede ser lento y existen fases de
estancamiento, estas mecánicas resultan especialmente útiles para mantener la
motivación.


\section{Motivación}

La escalada es un deporte eminentemente individual en el que el seguimiento del
progreso suele realizarse de forma manual y poco sistemática (libretas, memoria o
aplicaciones externas no integradas en el propio rocódromo). La motivación principal
de este proyecto es digitalizar dicha experiencia mediante una herramienta
centralizada, pensada para el contexto real de uso en rocódromo.

La aplicación propuesta permite registrar rutas completadas y conservar un historial
de actividad que facilite analizar la evolución del usuario. Además, incorpora
elementos de gamificación (puntos, logros, clasificaciones y retos) para transformar
el entrenamiento en una experiencia más participativa, social y motivadora,
favoreciendo la constancia y una competitividad saludable.


\section{Objetivos}

\subsection{Objetivo general}

Desarrollar una aplicación compatible con dispositivos móviles que permita a los
usuarios de un rocódromo registrar su actividad de escalada (bloque y vía) y que,
mediante técnicas de gamificación, fomente la motivación y la constancia en el
entrenamiento.

\subsection{Objetivos específicos}

\begin{enumerate}
	\item Implementar un sistema de gestión de usuarios que cubra registro,
	autenticación y perfil personal con consulta de rutas y estadísticas.
	\item Desarrollar un módulo de registro de escaladas que permita marcar rutas como
	\textit{flash}, completadas o proyecto.
	\item Incorporar mecánicas de gamificación orientadas a mejorar la motivación del
	usuario.
	\item Implementar funcionalidades sociales (amistades, comparación de progreso y
	estadísticas) para fomentar comunidad y competencia amistosa.
	\item Diseñar una interfaz y una experiencia de usuario claras, intuitivas y
	eficientes, adecuadas a escaladores con distintos niveles tecnológicos.
\end{enumerate}


\section{Estructura del documento}

La memoria se organiza de la siguiente forma. En el Capítulo~\ref{cap:introduccion}
se presenta el contexto, la motivación y los objetivos del proyecto. En el capítulo de
estado de la cuestión se revisan antecedentes y trabajos relacionados. A continuación,
el capítulo de descripción del trabajo detalla el análisis, diseño e implementación de
la solución propuesta. Posteriormente, el capítulo de conclusiones recoge los
resultados alcanzados y las líneas de trabajo futuro. Finalmente, se incluyen las
contribuciones personales, la bibliografía y los apéndices con material complementario.